% TEMPLATE-MILC-v3.tex - plantilla maestra UNICA de las guias MILC v3.
%
% La usan las tres familias de guias, cada una con su driver:
%   · Grados 6-11  -> scripts/build-guias-g11.py
%   · Bebras       -> scripts/build-guias-bebras.py
%   · Semillero    -> scripts/build-guias-semillero.py
%
% Antes habia tres copias casi identicas (~400 lineas iguales, 6 distintas):
% cada mejora habia que aplicarla tres veces, y en la practica no se hacia
% (el motor de recursos vivia solo en dos de ellas). Lo que varia entre
% familias esta parametrizado en 6 placeholders que cada driver llena
% (se nombran aqui SIN su sintaxis real: el builder reemplaza por texto plano,
% tambien dentro de los comentarios, y un valor multilinea rompe el preambulo):
%
%   HEADER_IZQ        encabezado izquierdo de pagina
%   HEADER_DER        encabezado derecho de pagina
%   PORTADA_ETIQUETA  rotulo sobre el numero grande (GRADO/MOMENTO/MODULO)
%   PORTADA_SERIE     linea de serie de la portada
%   PORTADA_ID        identificador de la guia en la portada
%   PORTADA_CIRCULO   el circulito de grado (°), o vacio si no aplica
%
% El resto de claves las documenta content/guias/_SCHEMA.md. El script valida
% que no queden placeholders sin reemplazar antes de escribir el archivo final.

\documentclass[11pt,letterpaper]{article}
\usepackage[margin=2.0cm]{geometry}
\usepackage[table]{xcolor}
\usepackage{fontspec}
\usepackage[spanish]{babel}
\usepackage{tikz}
% Librerías TikZ para los diagramas de las guías (bloque `recursos:`):
%  · babel  → babel en español vuelve activos caracteres como «>», y TikZ los usa
%             en sus opciones (>=stealth, ->). Sin esta librería, cualquier
%             diagrama con flechas revienta con «\language@active@arg> has an
%             extra }». Debe ir ANTES de que se dibuje nada.
%  · positioning        → sintaxis «below left=2pt and 3pt».
%  · decorations.pathreplacing → llaves ({brace}) para señalar rangos.
\usetikzlibrary{babel,positioning,decorations.pathreplacing}
\usepackage{graphicx}
% Circuitos: el trabajo de electrónica se hace hoy con micro:bit y sus sensores
% integrados (MakeCode, sin cableado), así que no se necesitan esquemáticos.
% Si algún día entran montajes con protoboard, basta descomentar esta línea:
% los diagramas .tex/.tikz declarados en `recursos:` ya se incrustan solos.
% \usepackage{circuitikz}
\usepackage[most]{tcolorbox}
\usepackage{tabularx}
\usepackage{array}
\usepackage{fancyhdr}
\usepackage{titlesec}
\usepackage[document]{ragged2e}  % bandera a la derecha en todo el cuerpo
\usepackage{enumitem}
\usepackage{microtype}

% Helvetica Neue no trae la flecha U+2192, asi que cada «→» escrito literalmente
% en el YAML se compilaba a NADA, en silencio (462 flechas en 61 guias: rutas del
% tipo «paleta → tipografia → lockup» salian rotas). Mapeamos el caracter a su
% version matematica, que si tiene glifo. Asi el autor puede seguir escribiendo
% «→» comodamente en el YAML.
\usepackage{newunicodechar}
\newunicodechar{→}{\ensuremath{\rightarrow}}
% Mismo problema con los simbolos de marca y vinetas: el «✓» de «una palomita ✓»
% salia como cuadrito vacio en el PDF de todas las guias que lo usaban.
\usepackage{pifont}
\newunicodechar{✓}{\ding{51}}
\newunicodechar{✗}{\ding{55}}
\newunicodechar{✔}{\ding{52}}
\newunicodechar{•}{\textbullet}
\newunicodechar{≥}{\ensuremath{\geq}}
\newunicodechar{≤}{\ensuremath{\leq}}

\setmainfont{Helvetica Neue}
\setsansfont{Helvetica Neue}
\setlength{\parindent}{0pt}
\setlength{\parskip}{6pt}
% Sin particion de palabras: en 6.º-7.º una palabra cortada por guion al final
% de linea («Comuni-cacion») frena la lectura mucho mas que una linea corta.
\hyphenpenalty=10000
\exhyphenpenalty=10000
\pretolerance=2000
\tolerance=3000
\setlength{\emergencystretch}{3em}
\linespread{1.10}
\renewcommand{\arraystretch}{1.25}
\setlist{leftmargin=5mm,itemsep=1mm,topsep=1mm}
\newcolumntype{Y}{>{\RaggedRight\arraybackslash}X}

\definecolor{milcMagenta}{HTML}{E6007E}
\definecolor{milcVino}{HTML}{5A0038}
\definecolor{milcTurquesa}{HTML}{0093A5}
\definecolor{milcVerde}{HTML}{58A923}
\definecolor{milcMostaza}{HTML}{E5B400}
\definecolor{milcOcre}{HTML}{B86600}
\definecolor{milcNegro}{HTML}{241816}
\definecolor{milcGris}{HTML}{F7F7F4}
\definecolor{milcLinea}{HTML}{E8E2DD}
\definecolor{milcGradePrimary}{HTML}{0D9488}
\definecolor{milcGradeDark}{HTML}{115E59}
\definecolor{milcGradeSoft}{HTML}{CCFBF1}
\definecolor{milcGradeAccent}{HTML}{CFE7FF}
\definecolor{milcGradeText}{HTML}{FFFFFF}
\color{milcNegro}

\pagestyle{fancy}
\fancyhf{}
\lhead{\small Territorio interior · Educación socioemocional}
\rhead{\small Grado 9° · Momento 9 · MILC v3}
\cfoot{\small\thepage}
\renewcommand{\headrulewidth}{0pt}

\newtcolorbox{titlebox}[1]{
  enhanced,colback=#1,colframe=#1,arc=7mm,boxrule=0pt,
  left=7mm,right=7mm,top=3mm,bottom=3mm,
  before skip=8pt,after skip=8pt,
  fontupper=\sffamily\bfseries\Large\color{white}
}

\newtcolorbox{softbox}[3]{
  enhanced,colback=#2,colframe=#1,borderline west={4pt}{0pt}{#1},
  arc=2mm,boxrule=.4pt,left=4mm,right=4mm,top=3mm,bottom=3mm,
  before skip=6pt,after skip=8pt,title={#3},coltitle=white,
  colbacktitle=#1,fonttitle=\sffamily\bfseries,fontupper=\RaggedRight,
  attach boxed title to top left={xshift=3mm,yshift=-2mm},
  boxed title style={arc=2mm, boxrule=0pt}
}

\newtcolorbox{drawbox}[2]{
  enhanced,colback=white,colframe=#1,arc=4mm,boxrule=1pt,
  left=5mm,right=5mm,top=4mm,bottom=4mm,before skip=8pt,after skip=8pt,
  title={#2},coltitle=white,colbacktitle=#1,fonttitle=\sffamily\bfseries,
  fontupper=\RaggedRight,
  attach boxed title to top left={xshift=4mm,yshift=-2mm},
  boxed title style={arc=3mm, boxrule=0pt}
}

\newtcolorbox{infoband}[1]{
  enhanced,colback=#1!8,colframe=#1!50,arc=2mm,boxrule=.5pt,
  left=4mm,right=4mm,top=2mm,bottom=2mm,before skip=4pt,after skip=4pt,
  fontupper=\small\RaggedRight
}

% Puente narrativo entre secciones (contrato editorial Conéctate).
% Da continuidad: "Acabas de X. Ahora vas a Y porque Z."
% Visualmente: banda izquierda en color del grado, texto en itálica pequeña.
\newtcolorbox{puentebox}{
  enhanced,colback=milcGradeSoft,colframe=milcGradeDark,
  borderline west={3pt}{0pt}{milcGradeDark},
  arc=0mm,boxrule=0pt,left=5mm,right=4mm,top=2.5mm,bottom=2.5mm,
  before skip=4pt,after skip=8pt,
  fontupper=\itshape\small\color{milcNegro}\RaggedRight
}
% La flecha va en modo matematico a proposito: Helvetica Neue NO tiene el glifo
% U+2192, asi que \textrightarrow se compilaba a NADA («Missing character: There
% is no → in font Helvetica Neue») y el puente salia sin flecha. En modo
% matematico la flecha viene de la fuente matematica y si se dibuja.
\newcommand{\puente}[1]{%
  \begin{puentebox}%
    \textcolor{milcGradeDark}{$\rightarrow$}\hspace{2mm}#1%
  \end{puentebox}%
}

\newcommand{\linead}[1]{\noindent\color{milcLinea}\rule{#1}{0.5pt}\color{milcNegro}}
\newcommand{\respuesta}[1]{\par\vspace{2mm}\foreach \n in {1,...,#1}{\noindent\color{milcLinea}\rule{\linewidth}{0.5pt}\par\vspace{5mm}}\color{milcNegro}}
\newcommand{\checkbox}{\textcolor{milcGradeDark}{\fbox{\phantom{x}}}\hspace{5pt}}

% ─── Listas aireadas para las guías socioemocionales ────────────────────
% Componer las verificaciones como «\checkbox texto\\» deja el texto que salta
% de línea debajo del cuadrito y apelmaza el bloque. Estos entornos alinean la
% continuación y airean los ítems. Los usa Territorio Interior; cualquier
% familia puede hacerlo.
\newlist{checklista}{itemize}{1}
\setlist[checklista]{
  label=\textcolor{milcGradeDark}{\fbox{\phantom{x}}},
  leftmargin=9mm, labelsep=3.2mm, itemsep=2.4mm,
  topsep=2.5mm, parsep=0pt, partopsep=0pt, before=\RaggedRight
}
\newlist{pasoslista}{enumerate}{1}
\setlist[pasoslista]{
  label=\textcolor{milcGradeDark}{\sffamily\bfseries\arabic*.},
  leftmargin=8mm, labelsep=3mm, itemsep=2.8mm,
  topsep=1mm, parsep=0pt, partopsep=0pt, before=\RaggedRight
}

\newcommand{\phasecard}[4]{
\begin{tcolorbox}[enhanced,colback=#3,colframe=#2,arc=3mm,boxrule=.5pt,height=3.05cm,valign=center,halign=center,left=1.5mm,right=1.5mm,top=2mm,bottom=2mm]
{\sffamily\bfseries\fontsize{22}{24}\selectfont\textcolor{#2}{#1}}\par
{\sffamily\bfseries\small #4}
\end{tcolorbox}
}

% ── Piezas del rediseno 2026 (legibilidad 6.º-11.º) ───────────────────
% Banda de estacion: reemplaza el titlebox macizo. Titulo a la izquierda,
% dato practico (tiempo/modalidad) a la derecha, en una sola linea.
\newcommand{\milcestacion}[2]{%
\begin{tcolorbox}[enhanced,colback=#1,colframe=#1,arc=2mm,boxrule=0pt,
  left=5mm,right=5mm,top=2.6mm,bottom=2.6mm,before skip=11pt,after skip=7pt]
{\sffamily\bfseries\fontsize{15}{18}\selectfont\color{white}#2}
\end{tcolorbox}}

% Tarjeta de paso numerada: sustituye las tablas «Pilar / Aplicacion», que
% eran formato de consulta docente, no de estudiante. El numero va en un
% circulo sobre el borde para que la secuencia se lea de un vistazo.
\newcommand{\milcpaso}[3]{%
\begin{tcolorbox}[enhanced,colback=white,colframe=milcLinea,arc=1.5mm,boxrule=.9pt,
  left=14mm,right=4mm,top=3mm,bottom=3mm,before skip=4pt,after skip=4pt,
  fontupper=\RaggedRight,
  overlay={\node[anchor=north west,circle,fill=#2,text=white,inner sep=0pt,
    minimum size=7.4mm,font=\sffamily\bfseries\small]
    at ([xshift=3.6mm,yshift=-3.2mm]frame.north west) {#1};}]
#3
\end{tcolorbox}}

% Casilla marcable de tamano util con lapiz (la anterior era \fbox{\phantom{x}},
% de alto variable segun el texto que la seguia).
\newcommand{\milcmarca}{\framebox[3.8mm]{\rule{0pt}{3.4mm}}}

% Fila marcable de lista suelta (checks de la estacion 1).
\newcommand{\milccheckrow}[1]{%
\par\vspace{1.8mm}\noindent
\makebox[6.4mm][l]{\raisebox{-.55mm}{\textcolor{milcNegro}{\milcmarca}}}%
\parbox[t]{\dimexpr\linewidth-6.4mm\relax}{\RaggedRight #1}\par}

% Celda marcable dentro de la rubrica (tabla de autoevaluacion). Misma
% sangria francesa, pero pensada para vivir dentro de una columna X.
\newcommand{\milcopcion}[1]{%
\makebox[5.2mm][l]{\raisebox{-.55mm}{\textcolor{milcNegro}{\milcmarca}}}%
\parbox[t]{\dimexpr\linewidth-5.2mm\relax}{\RaggedRight #1}}

% ── Recursos visuales de la guía ──────────────────────────────────────
% Declarados en el bloque `recursos:` del YAML y emitidos por el builder.
% \guiaFigura   → imagen rasterizada o PDF (foto, carta celeste, montaje).
% \guiaDiagrama → fuente TikZ/circuitikz (.tex) incrustada como vector:
%                 es la vía para circuitos y esquemáticos editables.
% #1 = ruta relativa al .tex · #2 = pie (puede ir vacío)
\newcommand{\guiaFigura}[2]{%
\begin{center}
\includegraphics[width=0.88\linewidth,height=0.38\textheight,keepaspectratio]{#1}\\[1.5mm]
{\footnotesize\itshape\textcolor{milcNegro!75}{#2}}
\end{center}
}

\newcommand{\guiaDiagrama}[2]{%
\begin{center}
\input{#1}\\[1.5mm]
{\footnotesize\itshape\textcolor{milcNegro!75}{#2}}
\end{center}
}

\begin{document}
\thispagestyle{empty}

% ── Portada ───────────────────────────────────────────────────────────
% Antes: un rectangulo de color a pagina completa. Se veia bien en pantalla,
% pero en la fotocopiadora del colegio se comia el toner y salia gris sucio.
% Ahora la banda ocupa el tercio superior y el resto va en blanco.
\begin{tikzpicture}[remember picture,overlay]
  \path[fill=milcGradePrimary]
    (current page.north west) rectangle ([yshift=-10.6cm]current page.north east);

  \node[anchor=north west,text=milcGradeText] at ([xshift=2.0cm,yshift=-1.55cm]current page.north west)
    {\sffamily\bfseries\fontsize{12}{14}\selectfont MOMENTO};
  \node[anchor=north west,text=milcGradeText,opacity=.85] at ([xshift=2.0cm,yshift=-2.35cm]current page.north west)
    {\sffamily\bfseries\fontsize{8.5}{10}\selectfont TERRITORIO INTERIOR · SOCIOEMOCIONAL};
  \node[anchor=north east,text=milcGradeText,opacity=.85,align=right] at ([xshift=-2.0cm,yshift=-1.55cm]current page.north east)
    {\sffamily\bfseries\fontsize{9}{12}\selectfont I.E. Sor María Juliana\\Cartago, Valle del Cauca};

  \node[anchor=west,text=milcGradeText] (gradonum) at ([xshift=1.85cm,yshift=-7.1cm]current page.north west)
    {\sffamily\bfseries\fontsize{112}{112}\selectfont 9};
  % El circulito de grado (°) solo aplica a las guias de grado («11°»); en
  % Bebras y Semillero el numero grande es un momento/modulo, no un grado.
  % Se posiciona relativo al nodo `gradonum`, no en coordenadas absolutas.
  

  \node[anchor=west,text=milcGradeText,text width=11.4cm,align=left] at ([xshift=1.5cm,yshift=0.5cm]gradonum.east)
    {
     {\sffamily\bfseries\fontsize{9}{11}\selectfont\MakeUppercase{Territorio interior · Socioemocional}}\\[3mm]
     {\sffamily\bfseries\fontsize{21}{25}\selectfont\setlength{\baselineskip}{27pt}Justicia\\[4pt]y venganza\\[4pt]reparar no es castigar}\\[3.5mm]
     {\sffamily\bfseries\fontsize{15}{18}\selectfont Grado 9° · Momento 9}
    };
\end{tikzpicture}

\vspace*{9.4cm}
\vfill

{\sffamily\bfseries\fontsize{8}{10}\selectfont\textcolor{milcGradeDark}{LO QUE TE LLEVAS HOY}}

\begin{tcolorbox}[enhanced,colback=white,colframe=milcNegro,arc=2mm,boxrule=1.6pt,
  left=5mm,right=5mm,top=4mm,bottom=4mm,before skip=3pt,after skip=12pt,
  fontupper=\RaggedRight\fontsize{11}{15.5}\selectfont]
Tu comparación de dos justicias: un mismo caso resuelto por la vía del castigo y por la vía de la reparación, con lo que gana y pierde cada parte en cada una, y las tres condiciones sin las cuales una reparación no se puede intentar.
\end{tcolorbox}

\begin{tcolorbox}[enhanced,colback=milcGris,colframe=milcGris,arc=2mm,boxrule=0pt,
  left=5mm,right=5mm,top=3.5mm,bottom=3.5mm,after skip=12pt]
{\sffamily\bfseries\fontsize{8}{10}\selectfont\textcolor{milcNegro!55}{ESTA GUÍA ES DE}}\\[3mm]
\begin{tabularx}{\linewidth}{X p{3.2cm} p{3.2cm}}
\linead{\linewidth} & \linead{3.0cm} & \linead{3.0cm}\\[-1mm]
{\fontsize{8.5}{10}\selectfont\textcolor{milcNegro!55}{Nombre completo}} &
{\fontsize{8.5}{10}\selectfont\textcolor{milcNegro!55}{Grupo}} &
{\fontsize{8.5}{10}\selectfont\textcolor{milcNegro!55}{Fecha}}\\
\end{tabularx}
\end{tcolorbox}

\vfill

\noindent\textcolor{milcLinea}{\rule{\linewidth}{1pt}}\\[2mm]
{\sffamily\bfseries\fontsize{9.5}{12}\selectfont Autor: Dr. Álvaro Cárdenas Orozco}\\[1mm]
{\sffamily\fontsize{8.5}{11}\selectfont\textcolor{milcNegro!55}{Metodología MILC · Apertura ancestral · Justicia restaurativa · Triángulo Dussel-Estoicismo-Floridi}}

\newpage

\section*{Apertura: Justicia y venganza: reparar y castigar no son lo mismo}

\begin{softbox}{milcGradeDark}{milcGradeSoft}{Saber ancestral}
Entre el pueblo \textbf{wayuu}, en La Guajira, existe un sistema normativo propio cuyo aplicador es el \textbf{pütchipü'üi} ---el \textbf{palabrero}---, y que la \textbf{UNESCO} inscribió en \textbf{2010} en la Lista Representativa del Patrimonio Cultural Inmaterial de la Humanidad. \emph{Ya conociste al palabrero en el momento 5 del grado 7, por su oficio: hablar en vez de responder con los puños.} \textbf{Hoy nos importa otra cosa: \emph{cómo termina} lo que él resuelve.} \textbf{El sistema wayuu está inspirado en la \emph{reparación} y la \emph{compensación}.} \emph{Es decir: cuando una familia ofende a otra, lo que se busca no es principalmente que el ofensor sufra ---es que \textbf{la familia ofendida quede compensada}---.} \textbf{Y eso cambia todo el procedimiento:} \emph{no hay que probar quién es peor persona, hay que establecer \textbf{qué se debe y cómo se paga}}. \textbf{Fíjate en la consecuencia más rara para nuestros oídos:} \emph{el acuerdo puede cerrarse \textbf{sin que nadie vaya preso}} ---y sin embargo el conflicto queda cerrado, y las dos familias pueden volver a convivir---. \textbf{Guarda la distinción:} \emph{una justicia le pregunta al daño «¿quién lo hizo y cuánto merece?»; la otra le pregunta «¿qué hace falta para arreglarlo?».}
\end{softbox}

\begin{softbox}{milcTurquesa}{milcGris}{Saber contemporáneo}
¿Y funciona? \textbf{Se midió, con el método más exigente que hay.} La \textbf{Colaboración Campbell} publicó en \textbf{2013} una revisión sistemática dirigida por \textbf{Heather Strang} y \textbf{Lawrence Sherman}, con \textbf{diez estudios} en los que víctimas y ofensores ---que habían aceptado encontrarse--- se asignaban \textbf{al azar} a un encuentro cara a cara de \textbf{justicia restaurativa} o a la justicia penal común. \textbf{Los resultados van en dos direcciones y las dos importan.} \textbf{Sobre los ofensores:} quienes pasaron por el encuentro \textbf{cometieron significativamente menos delitos después} que los del grupo de comparación ---y el efecto fue \textbf{mayor en los delitos violentos} que en los de propiedad---. \textbf{Sobre las víctimas, que es lo que casi nunca se mide:} quienes participaron quedaron \textbf{más satisfechas} con cómo se manejó su caso, \textbf{recibieron con más frecuencia una disculpa} y \textbf{la calificaron como sincera}, \textbf{quedaron \emph{menos} inclinadas a querer venganza} y \textbf{sufrieron menos síntomas de estrés postraumático} (Strang y otros, 2013). \emph{Léelo otra vez, porque desarma la intuición más común:} \textbf{no es que reparar sea más blando y por eso menos eficaz. Salió mejor en las dos columnas ---la del que hizo el daño y la del que lo recibió---.}
\end{softbox}

\begin{softbox}{milcMostaza}{milcGris}{Pregunta-puente}
\textit{El sistema wayuu pregunta \textbf{qué se debe}, no \textbf{cuánto merece sufrir}. \textbf{La última vez que algo te pareció injusto, ¿qué querías: que arreglaran el daño} \ldots \textbf{o que a alguien le fuera mal}?}
\end{softbox}

\begin{softbox}{milcVerde}{milcGris}{Saber y saber hacer}
Al terminar podrás: (1) \textbf{distinguir} las dos preguntas que puede hacerle una justicia a un daño; (2) \textbf{citar} lo que se midió sobre la justicia restaurativa, tanto en reincidencia como en lo que les pasó a las víctimas; (3) \textbf{resolver} un mismo caso por las dos vías y comparar qué gana y pierde cada parte; (4) \textbf{nombrar} las tres condiciones sin las cuales una reparación no debe siquiera intentarse.
\end{softbox}



\small
\begin{tabularx}{\textwidth}{>{\bfseries}p{3.0cm}Y}
\rowcolor{milcGradePrimary!12}Autor & Dr. Álvaro Cárdenas Orozco\\
DBA & Comprende los fenómenos que afectan a su territorio, regula sus emociones ante ellos y acompaña a otros con empatía (transversal 6°--11°, articulado con la política «Escuela, territorio de vida» del MEN, Resolución 006519 de 2025, y la Ley 1620 de 2013)\\
Referentes & Primera ayuda psicológica (OMS, 2011/2012) · Directrices IASC (2007) · USGS y Servicio Geológico Colombiano · Pueblo nasa y la minga (CRIC) · Dussel · Estoicismo · Floridi\\
Producto final & Tu comparación de dos justicias: un mismo caso resuelto por la vía del castigo y por la vía de la reparación, con lo que gana y pierde cada parte en cada una, y las tres condiciones sin las cuales una reparación no se puede intentar.\\
\end{tabularx}
\normalsize

\puente{En el momento 8 repararon un daño colectivo. \textbf{Hoy la pregunta se vuelve general:} \emph{cuando alguien hace daño, ¿para qué sirve lo que hacemos después?}}

\milcestacion{milcGradeDark}{Ruta MILC: así vamos a aprender}
\begin{tabularx}{\textwidth}{>{\centering\arraybackslash}X>{\centering\arraybackslash}X>{\centering\arraybackslash}X>{\centering\arraybackslash}X}
\phasecard{1}{milcVerde}{milcGris}{Escucha\\[-1mm]\small Observo, escucho y pregunto.} &
\phasecard{2}{milcTurquesa}{milcGris}{Sistematización\\[-1mm]\small Distingo reparar de cobrar.} &
\phasecard{3}{milcMagenta}{milcGris}{Praxis\\[-1mm]\small Construyo y aplico.} &
\phasecard{4}{milcVino}{milcGris}{Evaluación\\[-1mm]\small Reflexión liberadora.}
\end{tabularx}


\puente{Primero, tu propia intuición sin filtro: qué te parece que merece alguien que hizo algo malo. \emph{Escríbelo antes de leer nada.}}

\milcestacion{milcVerde}{Estación 1 de 3 · Escucha}

\begin{softbox}{milcVerde}{milcGris}{Escena para escuchar}
\textbf{Actividad 1 · IDENTIFICA --- Qué merece.} \textbf{Nadie va a leer esta página, y por eso hay que contestarla sin filtro.} \textbf{Primera parte.} Piensa en \textbf{algo injusto que te hayan hecho a ti} ---no lo peor: algo real---. Escribe \textbf{qué querías que pasara} en ese momento. \emph{Sin corregirte y sin quedar bien: si querías que a esa persona le fuera mal, escríbelo.} \textbf{Segunda parte.} Ahora responde tres preguntas sobre lo mismo: \emph{¿qué habría hecho falta para que el daño quedara \textbf{arreglado}? ¿Es lo mismo que querías en el momento? ¿Alguna de las dos cosas pasó?} \textbf{Tercera parte.} Piensa en algo que \textbf{tú} hiciste mal: \emph{¿qué te pasó después ---castigo, nada, tener que arreglarlo--- y cuál de esas cosas hizo que no lo repitieras?} \textbf{Y la última:} \emph{¿alguna vez te castigaron por algo y saliste con más rabia que arrepentimiento?}
\end{softbox}

{\sffamily\bfseries\fontsize{8}{10}\selectfont\textcolor{milcNegro!55}{MARCA CUANDO LO HAYAS HECHO}}
\milccheckrow{Escribiste qué querías en el momento, sin corregirte.}
\milccheckrow{Distinguiste eso de \textbf{qué habría hecho falta para arreglarlo}.}
\milccheckrow{Anotaste qué hizo que \textbf{tú} no repitieras algo.}

\begin{infoband}{milcVerde}


\textbf{En tu cuaderno (Actividad 1 · IDENTIFICA).} Título: «Actividad 1. Qué merece». \textbf{Formato:} qué querías + qué habría arreglado el daño + qué hizo que tú no repitieras algo + la vez que saliste con más rabia. \textbf{Extensión:} media página. \textbf{Sabes que terminaste cuando} puedes ver si \textbf{lo que querías} y \textbf{lo que habría arreglado el daño} eran la misma cosa. Esta página es tuya: \textbf{nadie más tiene que leerla si tú no quieres}.
\end{infoband}

\puente{Ya lo tienes. Ahora, las dos preguntas, lo que muestran los estudios y por qué la venganza no cierra nada.}

\milcestacion{milcTurquesa}{Fase 2 · Sistematización: entender las dos vías}

\begin{softbox}{milcTurquesa}{milcGris}{Cuatro claves sobre las dos justicias}
Cuatro claves para dejar de confundir que alguien pague con que algo se arregle.
\end{softbox}

\milcpaso{1}{milcTurquesa}{\textbf{Dos preguntas distintas frente al mismo daño.} \emph{La primera:} \textbf{«¿quién lo hizo y cuánto merece?»} ---mira hacia atrás, busca un culpable y fija un castigo proporcional---. \emph{La segunda:} \textbf{«¿qué hace falta para arreglarlo?»} ---mira hacia adelante, busca qué se debe y a quién---. \textbf{El sistema wayuu está construido sobre la segunda:} \emph{está inspirado en la reparación y la compensación}, y por eso un conflicto puede cerrarse \textbf{sin que nadie vaya preso} y con las dos familias conviviendo otra vez. \emph{Y ojo con una confusión frecuente:} \textbf{la segunda vía \emph{no} es más blanda}. \emph{Muchas veces es más exigente ---porque al que hizo el daño le toca dar la cara, oír lo que causó y \textbf{hacer algo}, en vez de simplemente aguantar una sanción y esperar a que pase---.}}
\milcpaso{2}{milcTurquesa}{\textbf{Lo que se midió, y salió mejor en las dos columnas.} La revisión Campbell de \textbf{2013}, con \textbf{diez estudios} de asignación aleatoria, encontró que quienes pasaron por un encuentro restaurativo \textbf{delinquieron significativamente menos después}, con \textbf{mayor efecto en delitos violentos}; y que las víctimas quedaron \textbf{más satisfechas}, \textbf{recibieron disculpas que consideraron sinceras}, \textbf{quedaron menos inclinadas a la venganza} y \textbf{tuvieron menos síntomas de estrés postraumático} (Strang y otros, 2013). \emph{Ese último dato merece un segundo:} \textbf{encontrarse con quien te hizo daño ---voluntariamente y bien acompañado--- dejó a las víctimas \emph{mejor}, no peor}. \textbf{Y fíjate en un detalle del diseño que importa mucho:} \emph{en esos estudios \textbf{las dos partes habían aceptado encontrarse antes} de saber en qué grupo caían}. \emph{Sin esa aceptación, nada de esto aplica.}}
\milcpaso{3}{milcTurquesa}{\textbf{La venganza no cierra: sostiene.} \emph{La venganza se siente como justicia porque promete un final ---«que sienta lo que yo sentí»---.} \textbf{El problema es que produce exactamente el material para la siguiente ronda:} \emph{ahora el otro tiene su propio agravio, igual de real que el tuyo.} \textbf{Es lo que en un colegio se ve como «la vuelta»,} \emph{y es la razón por la que unos conflictos duran años y nadie recuerda cómo empezaron ---exactamente como la cerca del momento 8 del grado 8---.} \textbf{Y hay algo más:} \emph{la venganza le devuelve el protagonismo al que hizo el daño}. \textbf{Reparar hace lo contrario: pone en el centro a quien lo recibió y lo que necesita.}}
\milcpaso{4}{milcTurquesa}{\textbf{Las tres condiciones, y sin ellas no se intenta.} \emph{Esto no sirve para todo ni siempre, y decirlo es parte de enseñarlo bien.} \textbf{Primera: voluntariedad de las dos partes} ---en los estudios, ambas habían aceptado encontrarse; \emph{obligar a alguien a «reconciliarse» es hacerle un segundo daño}---. \textbf{Segunda: reconocimiento previo de los hechos} ---si el que hizo el daño lo niega, no hay nada que reparar y el encuentro se vuelve un juicio---. \textbf{Tercera: que no haya desequilibrio de poder ni riesgo} ---por eso \textbf{esto no aplica al control coercitivo del momento 3, ni a las situaciones tipo III}: ahí sentar a las dos partes a conversar pone a una en peligro---. \emph{Cuando falta alguna de las tres, la respuesta correcta es la ruta institucional, no la mediación.}}

\begin{drawbox}{milcTurquesa}{Las dos vías, paso a paso}
\begin{pasoslista}
\item \textbf{Castigo:} se establece quién fue → se fija una sanción proporcional → se cumple → \emph{el caso se cierra formalmente}.
\item \textbf{Reparación:} se reconoce el hecho → \textbf{quien recibió el daño dice qué necesita} → se acuerda qué se hace → se cumple → \emph{se declara cerrado}.
\item \textbf{Diferencia clave:} \emph{en la primera, el afectado no tiene que decir nada. En la segunda, es quien manda.}
\item \textbf{Y no se excluyen:} \emph{muchos sistemas hacen las dos ---sanción \textbf{y} reparación---.}
\item \textbf{Antes de intentar la segunda:} \emph{voluntariedad, reconocimiento y ausencia de riesgo.}
\end{pasoslista}
\end{drawbox}

\begin{softbox}{milcMostaza}{milcGris}{Errores comunes y su corrección}
\textbf{«Si no hay castigo no hubo justicia»:} confunde el método con el objetivo. \textbf{Llamar reparación a un castigo suave:} reparar es hacer algo por el afectado, no sufrir menos. \textbf{Obligar a reconciliarse:} sin voluntariedad es un segundo daño. \textbf{Mediar cuando hay desequilibrio o riesgo:} pone en peligro a una de las partes. \textbf{Pedirle a la víctima que perdone:} el perdón no se cobra, como aprendiste en el grado 7. \textbf{Creer que reparar es más fácil:} suele ser más exigente. \textbf{Olvidar la garantía:} sin algo que cambie, el acuerdo dura lo que dura la reunión.
\end{softbox}

\begin{infoband}{milcTurquesa}


\textbf{En tu cuaderno (Actividad 2 · EXPLICA).} Título: «Actividad 2. Las dos preguntas». \textbf{Formato:} escribe las dos preguntas y, debajo de cada una, \textbf{qué tipo de solución produce}. Después anota los cuatro resultados que encontró la revisión Campbell. \textbf{Extensión:} media página. \textbf{Sabes que terminaste cuando} puedes explicar por qué \textbf{reparar no es la versión blanda de castigar}.
\end{infoband}

\puente{Ya sabes que hay dos vías. Es hora de recorrerlas las dos con el mismo caso, que es la única forma de compararlas de verdad.}

\milcestacion{milcMagenta}{Fase 3 · Praxis: resolver el mismo caso dos veces}

\begin{softbox}{milcMagenta}{milcGris}{Cómo se compara castigo con reparación}
El ejercicio es simple y difícil: el mismo caso, resuelto dos veces, y una comparación honesta que incluya lo que cada vía no logra.
\end{softbox}

\milcpaso{1}{milcMagenta}{\textbf{El caso, por la vía del castigo (10 min, en grupos).} Tomen \textbf{un caso inventado} ---algo dañado, algo dicho, algo tomado sin permiso--- y resuélvanlo \textbf{como se hace normalmente}: establezcan quién fue y fijen la sanción. \textbf{Escriban qué pasa después:} \emph{¿el afectado recupera algo? ¿el que lo hizo entendió qué causó? ¿se vuelven a hablar?}}
\milcpaso{2}{milcMagenta}{\textbf{El mismo caso, reparando (12 min).} Ahora el mismo caso por la otra vía, en el orden correcto: \textbf{reconocimiento del hecho → \emph{el afectado dice qué necesita} → acuerdo de qué se hace → garantía}. \emph{Cuidado con la trampa que casi todos los grupos caen:} \textbf{si su solución termina siendo «que le pida perdón delante de todos y no lo vuelva a hacer», eso es un castigo con otra ropa}. \emph{Reparar es que el afectado quede efectivamente compensado ---devolver, arreglar, corregir en público lo que se dijo, hacer algo concreto---.}}
\milcpaso{3}{milcMagenta}{\textbf{Quién gana qué (8 min).} Hagan una tabla de dos columnas ---castigo y reparación--- y tres filas: \emph{qué gana el afectado, qué gana el que hizo el daño, qué gana el curso}. \textbf{Y agreguen una cuarta fila, la más honesta: \emph{qué pierde cada vía}.} \emph{Porque la reparación también pierde algo:} \textbf{es lenta, exige que las dos partes quieran, y no sirve cuando alguien niega los hechos.}}
\milcpaso{4}{milcMagenta}{\textbf{Las tres condiciones (7 min).} Escríbanlas y después \textbf{apliquen el filtro a tres casos}: uno de convivencia normal, uno de acoso repetido y uno de los que en el grado 8 clasificaron como \textbf{tipo III}. \emph{Van a ver rápido cuál pasa el filtro y cuáles no.} \textbf{Y esa es la conclusión más importante del momento:} \emph{saber \textbf{cuándo esto no se usa} es tan parte de la herramienta como saber usarla.}}

\begin{drawbox}{milcMagenta}{Antes de cerrar tu comparación}
\begin{checklista}
\item Resolvieron \textbf{el mismo caso} por las dos vías, no dos casos distintos.
\item En la vía reparadora, \textbf{el afectado dice qué necesita} antes de que se acuerde nada.
\item Su solución reparadora \textbf{no es un castigo con otra ropa}.
\item La tabla incluye la fila de \textbf{qué pierde} cada vía.
\item Escribieron las \textbf{tres condiciones} y las aplicaron a tres casos.
\item Identificaron con claridad \textbf{cuándo esto no se usa}.
\end{checklista}
\end{drawbox}

\begin{softbox}{milcMagenta}{milcGris}{Plantilla de trabajo}
\textbf{Las dos preguntas.} \emph{«¿Quién lo hizo y cuánto merece?» frente a «¿qué hace falta para arreglarlo?»} \\[1mm]
\textbf{El orden de la reparación.} \emph{Reconocer → \textbf{preguntarle al afectado} → acordar → garantía → declarar cerrado.} \\[1mm]
\textbf{El filtro.} \emph{«¿Las dos partes quieren? · ¿Reconoce los hechos? · ¿Hay riesgo o desequilibrio?»} ---si falla una, va a la ruta---.
\end{softbox}

\begin{infoband}{milcMagenta}


\textbf{En tu cuaderno (Actividad 3 · CREA).} Título: «Actividad 3. El mismo caso, dos veces». \textbf{Formato:} el caso resuelto por las dos vías + la tabla de cuatro filas + las tres condiciones aplicadas a tres casos. \textbf{Extensión:} una página. \textbf{Sabes que terminaste cuando} puedes decir \textbf{en qué caso concreto \emph{no} usarías la reparación}.
\end{infoband}

\puente{El producto es una comparación honesta. \emph{Incluye lo que la reparación \emph{no} logra ---porque hay cosas que no logra---.}}

\milcestacion{milcMostaza}{Producto de la sesión}
\textbf{Comparación de dos justicias: el mismo caso por ambas vías, qué gana y pierde cada parte, y las tres condiciones}.
\begin{softbox}{milcMostaza}{milcGris}{Criterios de calidad}
(1) El mismo caso está resuelto por las dos vías, no dos casos distintos. (2) En la vía reparadora, la necesidad del afectado se establece antes del acuerdo. (3) La solución reparadora produce compensación efectiva, no una sanción reformulada. (4) La comparación incluye lo que cada vía pierde, no solo lo que gana. (5) Las tres condiciones están aplicadas a casos concretos, incluyendo uno donde la reparación no procede.
\end{softbox}

\puente{Antes de cerrar, revisa la trampa más común: si tu «solución reparadora» consiste en un castigo más suave, no entendiste la diferencia.}

\milcestacion{milcVino}{Evaluación liberadora · cinco dimensiones}

\begin{softbox}{milcVino}{milcGris}{Sentido de la evaluación}
Evaluar no es solo revisar una nota. Es reconocer qué aprendí, qué cambió en mí, qué emociones manejé, cómo me relaciono con la ciudad y la comunidad, y qué le dejaré a quien viene detrás.
\end{softbox}

\textbf{Mi reflexión liberadora en cinco dimensiones.} Completa cada frase iniciada:

\small
\begin{tabularx}{\textwidth}{>{\bfseries\RaggedRight\arraybackslash}p{3.4cm}Y>{\RaggedRight\arraybackslash}p{4.6cm}}
\rowcolor{milcVino}\color{white}Dimensión & \color{white}\bfseries Mi respuesta (frame starter) & \color{white}\bfseries Algo que descubrí\\
1. Personal & Lo que más cambió en mí fue \linead{6.5cm} & \linead{4.2cm}\\
2. Emocional & La emoción más fuerte fue \linead{2.5cm} y la regulé \linead{3cm} & \linead{4.2cm}\\
3. Ciudadana & En democracia esto importa porque \linead{6cm} & \linead{4.2cm}\\
4. Local (Cartago) & En mi barrio sería útil para \linead{6.5cm} & \linead{4.2cm}\\
5. Intergeneracional & A mi abuela/abuelo le diría \linead{6.5cm} & \linead{4.2cm}\\
\end{tabularx}
\normalsize

\begin{infoband}{milcVino}
\textbf{Para la próxima sustentación.} Vuelve a esta tabla en seis meses. Verás que las respuestas cambian — no porque lo aprendido cambie, sino porque tú cambias. La evaluación liberadora es una conversación contigo a través del tiempo.
\end{infoband}


\puente{Cerramos con tres voces: el rostro que reclama un derecho, la demora como remedio de la ira, y el cuidado de lo compartido.}

\milcestacion{milcGradeDark}{Triángulo de pensamiento · cierre filosófico}

\begin{softbox}{milcVino}{milcGris}{Dussel · lente del nosotros}
\textit{«El otro se revela realmente como otro\ldots como el pobre, el oprimido; el que, a la vera del camino, fuera del sistema, muestra su rostro sufriente y sin embargo desafiante: «¡Tengo hambre!, ¡tengo derecho a comer!»»}

{\footnotesize\itshape\color{milcNegro!70}--- Enrique Dussel · Filosofía de la liberación (1977)}

\textbf{Fíjate en quién habla en la cita: \emph{el que fue dañado}.} \emph{No un juez, no un tercero ---él mismo, y reclamando.} \textbf{Esa es exactamente la diferencia entre las dos justicias.} \emph{En la vía del castigo, el afectado no tiene que decir nada:} el sistema establece los hechos y fija una pena, \textbf{y él termina siendo un testigo de su propio caso}. \emph{En la vía de la reparación, en cambio, es quien manda:} \textbf{nada se acuerda antes de que él diga qué necesita}. \emph{Y por eso los datos de las víctimas en la revisión Campbell no sorprenden tanto como parece:} \textbf{la gente queda mejor cuando le devuelven la voz sobre lo que le pasó.}

\textbf{Cuando se arregla algo en mi curso, ¿le preguntan al afectado qué necesita, o le informan la decisión?}
\end{softbox}
\linead{\linewidth}\\[1mm]
\linead{\linewidth}

\begin{softbox}{milcMostaza}{milcGris}{Estoicismo · Séneca · Sobre la ira, libro II (c. 45 d.C.) · cuidado interior}
\textit{«El mayor remedio para la ira es la demora.»}

{\footnotesize\itshape\color{milcNegro!70}--- Séneca · Sobre la ira, libro II (c. 45 d.C.)}

\textbf{La venganza es la ira convertida en plan, y por eso la demora es su antídoto exacto.} \emph{«La vuelta» casi nunca se da en el momento ---se organiza---,} \textbf{y en ese intervalo cabe la única pregunta que la desactiva:} \emph{«¿esto va a arreglar algo, o solo va a hacer que el otro tenga ahora su propio motivo?»}. \textbf{Y hay una segunda lectura, importante para no usar mal a Séneca:} \emph{la demora \textbf{no} significa aguantarse un daño}. \emph{Significa esperar lo suficiente para escoger la respuesta \textbf{en vez de reaccionar}} ---y una de las respuestas posibles, la más incómoda para el orgullo, es \textbf{pedir que se repare}---.

\textbf{La última vez que quise devolvérsela a alguien, ¿qué habría arreglado eso?}
\end{softbox}
\linead{\linewidth}\\[1mm]
\linead{\linewidth}

\begin{softbox}{milcTurquesa}{milcGris}{Floridi · lente de la infoesfera}
\textit{«El florecimiento de las entidades informacionales y de toda la infoesfera debe promoverse preservando, cultivando y enriqueciendo sus propiedades.»}

{\footnotesize\itshape\color{milcNegro!70}--- Luciano Floridi · La ética de la información (2013; trad.)}

\textbf{Hoy la venganza tiene una herramienta que no existía y que la vuelve mucho peor: se puede hacer en público y queda escrita.} \emph{Exponer a alguien, publicar una captura, armar un grupo en su contra ---eso se siente como justicia porque «que todos sepan lo que hizo»---.} \textbf{Pero mira lo que produce:} \emph{un daño que ya no se puede reparar, y que además va a estar ahí dentro de tres años, cuando esa persona ya sea otra}. \textbf{Y produce algo más, que es lo que le importa a Floridi:} \emph{ensucia el espacio de todos} ---el curso entero queda con eso circulando---. \textbf{Por eso el castigo público en línea es la única forma de justicia que \emph{garantiza} que el conflicto no se cierre.}

\textbf{¿He visto que «exponer» a alguien haya arreglado algo alguna vez?}
\end{softbox}
\linead{\linewidth}\\[1mm]
\linead{\linewidth}

\begin{infoband}{milcNegro}
\textbf{Nota para el docente.} Las tres son citas textuales y llevan su referencia debajo. Puedes remitir a la obra en clase.
\end{infoband}

\milcestacion{milcVino}{Mi autoevaluación · marca una casilla por fila}

{\small
\setlength{\extrarowheight}{2.2pt}
\begin{tabularx}{\textwidth}{>{\RaggedRight\arraybackslash\bfseries}p{3.0cm}YYY}
\rowcolor{milcVino}\color{white}\bfseries Criterio & \color{white}\bfseries Lo logré & \color{white}\bfseries En proceso & \color{white}\bfseries Necesito apoyo\\[.6mm]
Comprensión del tema & \milcopcion{Explico las ideas con mis palabras.} & \milcopcion{Reconozco las ideas, falta precisión.} & \milcopcion{Necesito repasar lo básico.}\\[.8mm]
\rowcolor{milcGris}Calidad de la comparación & \milcopcion{Resolví el mismo caso por las dos vías y puedo decir qué gana y pierde cada parte en cada una.} & \milcopcion{Entiendo la diferencia, pero mis soluciones «reparadoras» siguen siendo castigos suaves.} & \milcopcion{Sigo pensando que si no hay castigo no hubo justicia.}\\[.8mm]
Aplicación y producto & \milcopcion{Mi producto cumple los criterios.} & \milcopcion{Está armado, con ajustes pendientes.} & \milcopcion{Aún está en construcción.}\\[.8mm]
\rowcolor{milcGris}Reflexión 5 dimensiones & \milcopcion{Respondí cada una con honestidad.} & \milcopcion{Respondí algunas con detalle.} & \milcopcion{Mis respuestas son muy generales.}\\[.8mm]
Triángulo de pensamiento & \milcopcion{Conecté las 3 voces con mi trabajo.} & \milcopcion{Conecté al menos una.} & \milcopcion{No las relaciono con mi proceso.}\\[.8mm]
\end{tabularx}}

\vspace{3mm}
\textbf{Hoy entendí que:}\\[1mm]
\linead{\linewidth}\\[1mm]
\linead{\linewidth}

\textbf{Mi compromiso para la próxima sesión es:}\\[1mm]
\linead{\linewidth}\\[1mm]
\linead{\linewidth}

\begin{softbox}{milcMostaza}{milcGris}{Cita de despedida · Marco Aurelio}
\textit{``El obstáculo en el camino se vuelve el camino. Lo que estorba al camino se vuelve camino.''}\\[1mm]
Cada error de hoy, bien observado, se convierte en pieza del camino que pisarás mañana.
\end{softbox}

\small
\begin{tabularx}{\textwidth}{>{\bfseries}p{3.6cm}Y}
\rowcolor{milcGradeDark!12}Nombre completo: & \linead{10cm}\\
Fecha: & \linead{4cm} \quad Curso/grupo: \linead{4cm}\\
Mi firma: & \linead{9cm}\\
\end{tabularx}
\normalsize

\begin{infoband}{milcGradeDark}
\textbf{Para la próxima sesión.} Dos cosas. \textbf{Primero, aplica la pregunta reparadora una vez} ---en algo pequeño de esta semana---: cuando algo te moleste, en vez de pensar qué merece el otro, pregúntate \emph{qué haría falta para arreglarlo} \textbf{y pídelo}. \emph{Anota qué pasó, sobre todo si no funcionó: también eso es un dato.} \textbf{Segundo, pregunta en tu casa cómo se arreglaban los daños:} averigua con alguien mayor \textbf{qué pasaba en su pueblo o su barrio cuando alguien dañaba algo de otro} ---un animal, una cerca, una cosecha, una ventana---: \emph{¿se llamaba a la autoridad, se pagaba, se reponía, se arreglaba entre las familias?} \textbf{Trae:} tu intento y lo que te contaron. \emph{Vas a encontrar, con mucha frecuencia, la lógica wayuu funcionando sin que nadie la llame así: en el campo colombiano el daño casi siempre se \textbf{reponía} ---se devolvía el animal, se pagaba la cosecha, se arreglaba la cerca---, y solo cuando eso fallaba se iba donde la autoridad.}
\end{infoband}

\end{document}
